\section{Lattices, fields and averaging operators}\label{sec:setup}

This section fixes notation and collects the elementary identities used throughout. Nothing here is deep, but the identities of Lemma~\ref{lem:1d} and the restriction estimate of Lemma~\ref{lem:restriction} are used repeatedly.

\subsection{Lattices and blocks}

Fix integers $d\ge2$, $L\ge2$, $k\ge1$ and $n\ge8$, and put
\[
N=L^k,\qquad \eta=\frac1N .
\]
The \emph{unit torus} and the \emph{fine torus} are
\[
\T=(\Z/n\Z)^d,\qquad \Tf=(\eta\Z/n\Z)^d .
\]
Thus $\Tf$ is a torus of side $n$ with lattice spacing $\eta$, and $\T\subset\Tf$ is the sublattice of points with integer coordinates. We write $e_1,\dots,e_d$ for the unit coordinate vectors, and $|x-x'|$ for the Euclidean distance on the torus (the minimum over representatives). Since $n\ge8$, every subset of a cube of side at most $7$ can be identified with a subset of $\eta\Z^d$, and we do this tacitly. On cubes of side at most $4\le n/2$ the torus distance coincides with the distance in $\eta\Z^d$. On larger cubes it can only be smaller, which makes H\"older quotients computed with the torus distance larger; this direction is harmless wherever such quotients appear as hypotheses.

For $y\in\T$ the \emph{block} with corner $y$ is
\[
B(y)=\{x\in\Tf:\ y_\mu\le x_\mu<y_\mu+1,\ \mu=1,\dots,d\}.
\]
It contains $N^d$ fine points, the blocks partition $\Tf$, and we write $y(x)$ for the corner of the block containing $x$. The only role of $L$ is that $N=L^k$. Our constants never depend on $L$, and the reader may think of $N$ as an arbitrary positive integer.

\begin{remark}[Relation to Ba\l aban's notation]
In \cite{BalPropI} the fine lattice is $T_\eta$, the unit lattice is $T_1^{(k)}$, and the fields have been rescaled so that the averaging blocks have side $1$. Our $\Tf$ and $\T$ are these lattices, with the torus side denoted by $n$. All of Ba\l aban's operators after $k$ renormalization steps live on this pair of lattices.
\end{remark}

\begin{figure}[t]
\centering
\begin{tikzpicture}[scale=1.5]
% blocks
\fill[fblue!12] (0,0) rectangle (1,1);
\draw[forange,very thick] (-0.03,-0.03) rectangle (2.03,1.03);
\foreach \i in {0,1,2,3} \draw[gray!60] (\i,0)--(\i,2);
\foreach \j in {0,1,2} \draw[gray!60] (0,\j)--(3,\j);
\foreach \i in {0,...,11} \foreach \j in {0,...,7} \fill[gray] (\i*0.25,\j*0.25) circle (0.9pt);
\foreach \i in {0,1,2} \foreach \j in {0,1} \fill[black] (\i,\j) circle (1.8pt);
\node[below left] at (0,0) {$y$};
\node[below] at (1,0) {$y+e_1$};
\node[fblue,fill=fblue!12,inner sep=1pt] at (0.5,0.62) {$B(y)$};
% weight profile of Q_1
\draw[forange,thick] (0,2.35)--(1,2.85)--(2,2.35);
\draw[gray,->] (-0.1,2.35)--(2.3,2.35) node[right,black] {\small $x_1$};
\node[forange,above] at (1,2.85) {\small weight of $Q_1$};
\node[align=left,right] at (3.2,1) {\small fine points $\Tf$ (spacing $\eta$)\\[2pt] \small unit points $\T$ (large dots)\\[2pt] \small {\color{fblue}block $B(y)$}: support of $Q_0$ at $y$\\[2pt] \small {\color{forange}tube $B(y)\cup B(y+e_1)$}: support of $Q_1$ at $y$};
\end{tikzpicture}
\caption{The two lattices for $d=2$, $N=4$. The block average $(Q_0f)(y)$ is the mean of $f$ over $B(y)$. The tube average $(Q_1f)(y)$ is the mean over $B(y)$ of the averages of $f$ along unit segments in direction $e_1$. Its weight is the triangle shown above the tube.}
\label{fig:lattices}
\end{figure}

\subsection{Fields, derivatives and the Hodge identity}

A \emph{scalar field} is a function $f:\Tf\to\R$. A \emph{gauge field} (a $1$-form) is a $d$-tuple $A=(A_1,\dots,A_d)$ of scalar fields. Here $A_\mu(x)$ is the value of $A$ on the oriented bond $(x,x+\eta e_\mu)$, and the value on the reversed bond is $-A_\mu(x)$. A $2$-form is a collection $F=(F_{\mu\nu})$ with $F_{\nu\mu}=-F_{\mu\nu}$. Here $F_{\mu\nu}(x)$ is the value on the plaquette spanned by $x$, $\eta e_\mu$ and $\eta e_\nu$. The inner products are
\[
\ip fg=\eta^d\sum_{x\in\Tf}f(x)g(x),\qquad
\ip AB=\sum_{\mu}\ip{A_\mu}{B_\mu},\qquad
\ip FE=\sum_{\mu<\nu}\ip{F_{\mu\nu}}{E_{\mu\nu}},
\]
and $\|\cdot\|$ denotes the corresponding norms. For a set $X\subset\Tf$, $\|f\|_X$ is the norm of $\one_Xf$. Unit-lattice fields are defined in the same way on $\T$, with inner product $\sum_{y\in\T}$ and no factor $\eta^d$. With these normalizations a block has volume $\eta^dN^d=1$, so constant fields have comparable norms on both lattices.

The forward and backward derivatives are
\[
(\partial_\mu f)(x)=\frac{f(x+\eta e_\mu)-f(x)}{\eta},\qquad (\partial_\mu^*f)(x)=\frac{f(x-\eta e_\mu)-f(x)}{\eta}.
\]
Then $\partial_\mu^*$ is the adjoint of $\partial_\mu$. All $\partial_\mu$ and $\partial_\nu^*$ commute with one another and with translations of $\Tf$. The exterior derivative and its adjoint are
\begin{align*}
(df)_\mu&=\partial_\mu f, & d^*A&=\sum_\mu\partial_\mu^*A_\mu,\\
(dA)_{\mu\nu}&=\partial_\mu A_\nu-\partial_\nu A_\mu, & (d^*F)_\nu&=\sum_\mu\partial_\mu^*F_{\mu\nu}.
\end{align*}
One checks directly that $d^*$ is the adjoint of $d$ and that $d(df)=0$. The (non-positive) Laplacian is $\Lap=-\sum_\mu\partial_\mu^*\partial_\mu=-d^*d$ on scalar fields. It acts on gauge fields componentwise. We use the shorthand
\[
\|\partial f\|^2=\sum_\mu\|\partial_\mu f\|^2=\ip f{-\Lap f},\qquad
\|\partial A\|^2=\sum_{\mu,\nu}\|\partial_\nu A_\mu\|^2 .
\]
Gauge fields of the form $d\lambda$ are called \emph{exact}, and those with $d^*A=0$ \emph{co-closed}.

\begin{lemma}[Hodge identities]\label{lem:hodge}
For every gauge field $A$,
\[
(d^*d+dd^*)A=-\Lap A\quad\text{(componentwise)},\qquad
\|dA\|^2+\|d^*A\|^2=\|\partial A\|^2 .
\]
Moreover $-\Lap$ is invertible on scalar fields with mean zero, and every gauge field decomposes orthogonally as $A=d\lambda+A'$ with $d^*A'=0$ and $\lambda=(-\Lap)^{-1}d^*A$. The same statements hold on the unit torus, with $\partial_\mu$ replaced by the unit difference $D_\mu g(y)=g(y+e_\mu)-g(y)$.
\end{lemma}

\begin{proof}
For the first identity,
\[
(d^*dA)_\mu=\sum_\nu\partial^*_\nu(\partial_\nu A_\mu-\partial_\mu A_\nu),\qquad (dd^*A)_\mu=\partial_\mu\sum_\nu\partial^*_\nu A_\nu,
\]
and the mixed terms cancel because $\partial_\mu$ and $\partial^*_\nu$ commute. Taking the inner product with $A$ gives the second identity. Next, $-\Lap$ is symmetric and $\ip f{-\Lap f}=\|\partial f\|^2$. This vanishes only for constant $f$ on the connected torus. Hence $-\Lap$ is invertible on the orthogonal complement of the constants. Since $d^*A$ has mean zero ($\ip1{d^*A}=\ip{d1}A=0$), $\lambda$ is well defined. Then $d^*(A-d\lambda)=d^*A+\Lap\lambda=0$, and $\ip{d\lambda}{A'}=\ip\lambda{d^*A'}=0$. The proof on $\T$ is identical.
\end{proof}

On the unit torus we write $\dl$ for the exterior derivative built from $D_\mu$, and $\dl^*$ for its adjoint, with $(D_\mu^*g)(y)=g(y-e_\mu)-g(y)$.

\subsection{Directional convolutions}

The flat block and tube averages introduced in this section are obtained by restricting products of one-dimensional averages to the unit lattice, so we introduce a convenient language for those products.

\begin{definition}
Let $\mu\in\{1,\dots,d\}$ and let $\kappa:\eta\Z\to\R$ vanish outside an interval of length at most $6$. The \emph{directional convolution} in direction $\mu$ with kernel $\kappa$ is
\[
(Kf)(x)=\sum_{s\in\eta\Z}\eta\,\kappa(s)\,f(x+se_\mu),
\]
where $x+se_\mu$ is taken modulo the torus. Products of directional convolutions in different directions are called \emph{product convolutions}.
\end{definition}

Directional convolutions commute with each other and with all translations of $\Tf$. The adjoint of $K$ has kernel $s\mapsto\kappa(-s)$. The composition of two convolutions in the same direction has kernel $(\kappa_1\star\kappa_2)(s)=\sum_t\eta\,\kappa_1(t)\kappa_2(s-t)$. Moreover $|\kappa_1\star\kappa_2|\le\|\kappa_1\|_\infty\sum_t\eta|\kappa_2(t)|$, and the support lengths add. The box averages $a_\mu$, $a^*_\mu$ and the triangle average $c_\mu$ have nonnegative kernels of mass $1$, bounded by $1$. Hence every composition of such averages, possibly composed with one further convolution whose kernel is bounded by $M$, has a kernel bounded by $1$ (respectively $M$). Its support length is the sum of the support lengths.

We use the following operators in direction $\mu$:
\begin{itemize}
\item the forward and backward \emph{box averages} $a_\mu$, $a_\mu^*$, with kernels $\one_{[0,1)}$ and $\one_{(-1,0]}$:
\[
(a_\mu f)(x)=\eta\sum_{j=0}^{N-1}f(x+j\eta e_\mu),\qquad (a_\mu^*f)(x)=\eta\sum_{j=0}^{N-1}f(x-j\eta e_\mu);
\]
\item the unit difference $D_\mu f(x)=f(x+e_\mu)-f(x)$ and the shift $(\tau_\mu f)(x)=f(x+(1-\eta)e_\mu)$;
\item the \emph{triangle average} $c_\mu=a_\mu a_\mu^*$.
\end{itemize}

\begin{lemma}[One-dimensional identities]\label{lem:1d}
For each $\mu$:
\begin{enumerate}[label=(\alph*)]
\item $D_\mu=a_\mu\partial_\mu=\partial_\mu a_\mu$ and $D_\mu^*=a_\mu^*\partial_\mu^*$;
\item $a_\mu=\tau_\mu a_\mu^*$;
\item $c_\mu$ has kernel $c(s)=(1-|s|)_+$, and $\sum_s\eta\,c(s)=1$;
\item $1-c_\mu=e_\mu\partial_\mu^*\partial_\mu$, where $e_\mu$ has a kernel $k$ with $0\le k\le\frac12$, $k(s)=0$ for $|s|\ge1$, and $|k(s)-k(s')|\le|s-s'|$. Consequently $e_\mu\partial_\mu$ is a directional convolution with kernel $k_1(s)=(k(s-\eta)-k(s))/\eta$, which satisfies $|k_1|\le1$ and vanishes outside $[-1,1+\eta]$;
\item $a_\mu-1=g_\mu\partial_\mu$, where $g_\mu$ has kernel $(1-\eta-s)\one_{[0,1)}(s)\in[0,1]$.
\end{enumerate}
\end{lemma}

\begin{proof}
All identities are between convolution operators, so it suffices to check them on $\eta\Z$, i.e.\ for functions of one variable $t$. The torus versions follow by periodization.

(a) $(a\partial f)(t)=\eta\sum_{j=0}^{N-1}\frac{f(t+(j+1)\eta)-f(t+j\eta)}{\eta}=f(t+1)-f(t)$, since $N\eta=1$. Taking adjoints gives the second identity.

(b) $(\tau a^*f)(t)=\eta\sum_{j=0}^{N-1}f(t+1-\eta-j\eta)=\eta\sum_{j'=0}^{N-1}f(t+j'\eta)$, with $j'=N-1-j$.

(c) The kernel of $aa^*$ at $s\in\eta\Z$ is $\eta\cdot\#\{t\in\eta\Z: t\in[0,1),\ s-t\in(-1,0]\}$. For $0\le s<1$ the admissible $t$ are $s,s+\eta,\dots,1-\eta$, which gives $\eta\cdot(1-s)/\eta=1-s$. The case $s<0$ is symmetric. Finally $\sum_s\eta\,c(s)=\eta\big(1+2\sum_{j=1}^{N-1}(1-j\eta)\big)=\eta(2N-1-(N-1))=1$.

(d) Let $-\Lap_1h(t)=\eta^{-2}(2h(t)-h(t+\eta)-h(t-\eta))$ denote the one-dimensional operator $\partial^*\partial$. A direct check shows $-\Lap_1(-|t|/2)=\eta^{-1}\delta_{t,0}$. Put
\[
k(t)=-\frac{|t|}2+\frac12\sum_s\eta\,|t-s|\,c(s).
\]
Then $-\Lap_1k=\eta^{-1}\delta_{\cdot,0}-c$. Hence the convolution $e$ with kernel $k$ satisfies $e\,\partial^*\partial f=(-\Lap_1k)\star f=f-cf$. If $|t|\ge1$, then $|t-s|=|t|-\operatorname{sgn}(t)s$ for all $s$ in the support of $c$. Since $c$ has mass $1$ and mean $0$, this gives $k(t)=0$. For $|t|<1$, the triangle inequality and $\sum\eta c(s)(t-s)=t$ give $k(t)\ge0$, while $\sum_s\eta|t-s|c(s)\le|t|+1$ gives $k(t)\le\frac12$. Each of the two terms defining $k$ is Lipschitz with constant $\frac12$. The formula for the kernel of $e\partial$ follows by summation by parts: $\sum_s\eta k(s)\frac{f(t+s+\eta)-f(t+s)}\eta=\sum_s\eta\frac{k(s-\eta)-k(s)}\eta f(t+s)$.

(e) $(a f)(t)-f(t)=\eta\sum_{j=0}^{N-1}\sum_{i=0}^{j-1}\eta\,\partial f(t+i\eta)=\sum_{i=0}^{N-1}\eta\,\eta(N-1-i)\,\partial f(t+i\eta)$, and $\eta(N-1-i)=1-\eta-i\eta$.
\end{proof}

\subsection{Averaging operators}

Let $\iota$ denote restriction to the unit torus, $(\iota f)(y)=f(y)$ for $y\in\T$. Its adjoint is $(\iota^*g)(x)=\eta^{-d}g(x)$ if $x\in\T$ and $0$ otherwise. Put $b=a_1a_2\cdots a_d$.

\begin{definition}[Averaging operators]\label{def:averages}
The \emph{block average} $Q_0$ and the \emph{tube averages} $Q_\mu$, $\mu=1,\dots,d$, map scalar fields on $\Tf$ to scalar fields on $\T$:
\begin{align*}
(Q_0f)(y)&=(\iota bf)(y)=\eta^d\sum_{x\in B(y)}f(x),\\
(Q_\mu f)(y)&=(\iota a_\mu bf)(y)=\eta^d\sum_{x\in B(y)}\ \eta\sum_{j=0}^{N-1}f(x+j\eta e_\mu).
\end{align*}
The \emph{gauge-field average} maps gauge fields on $\Tf$ to gauge fields on $\T$:
\[
(\bQ A)_\mu=Q_\mu A_\mu,\qquad\text{i.e.}\qquad (\bQ A)_\mu(y)=\eta^d\sum_{x\in B(y)}\eta\sum_{j=0}^{N-1}A_\mu(x+j\eta e_\mu).
\]
We write $Q_\sigma$, $\sigma\in\{0,1,\dots,d\}$, when a statement holds for all of them, and $\cK_\sigma$ for the product convolution with $Q_\sigma=\iota\cK_\sigma$ ($\cK_0=b$, $\cK_\mu=a_\mu b$).
\end{definition}

The gauge-field average is the linearization of Ba\l aban's averaging at $U=1$. In \cite[(1.11), (1.18)]{BalPropI} it is written as $(Q_kA)_b=\sum_{x\in B_k(b_-)}\eta^{d+1}A([x,x(b)])$, where $A([x,x(b)])$ is the sum of $A$ over the unit segment from $x$ to $x+e_\mu$. This is exactly $(\bQ A)_\mu(y)$. Ba\l aban's original averaging \cite[(1.8)]{BalPropI} also contains contour terms. These vanish in the axial gauge, and dropping them is the ``second point of view'' of \cite[(1.13)]{BalPropI}, which we adopt.

\begin{lemma}[Properties of the averages]\label{lem:avg}
\begin{enumerate}[label=(\alph*)]
\item $(Q_\sigma f)(y)=\eta^d\sum_xw_\sigma(y,x)f(x)$ with $0\le w_\sigma\le1$, $\eta^d\sum_xw_\sigma(y,x)=1$ and $\sum_yw_\sigma(y,x)=1$. Moreover $w_0(y,x)=\one_{B(y)}(x)$, and $w_\mu(y,\cdot)$ vanishes outside $B(y)\cup B(y+e_\mu)$.
\item $Q_\sigma1=1$ and $Q_\sigma^*1=1$. Further, $(Q_0^*g)(x)=g(y(x))$, $Q_0Q_0^*=I$, and $Q_0^*Q_0$ is the orthogonal projection onto fields that are constant on each block.
\item $|(Q_\sigma f)(y)|^2\le\eta^d\sum_xw_\sigma(y,x)|f(x)|^2$. Hence $\|Q_\sigma\|\le1$, $\|\bQ\|\le1$, and $|(Q_\sigma^* g)(x)|\le\max\{|g(y)|:w_\sigma(y,x)\neq0\}$, where at most two $y$ occur.
\item \emph{(Gauge covariance.)} $D_\mu Q_0=Q_\mu\partial_\mu$, i.e.\ $\bQ\,d=\dl Q_0$. More generally, $D_\mu Q_\sigma=\iota\,a_\mu\cK_\sigma\partial_\mu$.
\item \emph{(Curvature.)} $\dl\bQ=\bQ^{(2)}d$, where $(\bQ^{(2)}F)_{\mu\nu}=\iota\,a_\mu a_\nu b\,F_{\mu\nu}$.
\end{enumerate}
\end{lemma}

\begin{proof}
(a) For $Q_0$ this is the definition. For $Q_\mu$, the coefficient of $f(x)$ in $(Q_\mu f)(y)$ is $\eta^{d}\cdot\eta\,\#\{j\in\{0,\dots,N-1\}:x-j\eta e_\mu\in B(y)\}$. So $w_\mu(y,x)\le\eta N=1$, and $w_\mu(y,x)\neq0$ only if $x\in B(y)\cup B(y+e_\mu)$. The normalization $\eta^d\sum_xw_\sigma=Q_\sigma1=1$ is clear. The identity $\sum_yw_\sigma(y,x)=1$ is $Q_\sigma^*1=1$, which follows from $(b^*\iota^*g)(x)=g(y(x))$ and $a_\mu^*1=1$.

(b) This follows from (a). Moreover $(Q_0Q_0^*g)(y)=\eta^d\sum_{x\in B(y)}g(y)=g(y)$. Hence $Q_0^*Q_0$ is a self-adjoint idempotent, and its range consists of the block-constant fields.

(c) The first bound is Cauchy--Schwarz with the probability weights $\eta^dw_\sigma(y,\cdot)$. Summing over $y$ and using $\sum_yw_\sigma(y,x)=1$ gives $\|Q_\sigma f\|^2\le\|f\|^2$. (Below we sometimes use the weaker bound $\|Q_\sigma\|^2\le2$.) The bound on $Q_\sigma^*g$ follows from $(Q_\sigma^*g)(x)=\sum_yw_\sigma(y,x)g(y)$ and $\sum_yw_\sigma(y,x)=1$.

(d) $\iota$ commutes with unit translations, so $D_\mu Q_\sigma=\iota D_\mu\cK_\sigma=\iota\cK_\sigma a_\mu\partial_\mu$ by Lemma~\ref{lem:1d}(a). For $\sigma=0$ this equals $\iota a_\mu b\,\partial_\mu=Q_\mu\partial_\mu$.

(e) $(\dl\bQ A)_{\mu\nu}=D_\mu Q_\nu A_\nu-D_\nu Q_\mu A_\mu=\iota\,a_\nu b\,a_\mu\partial_\mu A_\nu-\iota\,a_\mu b\,a_\nu\partial_\nu A_\mu$, by (a) of Lemma~\ref{lem:1d}.
\end{proof}

Part (d) is the reason for the particular form of $\bQ$. A gauge transformation $A\mapsto A+d\lambda$ changes the average by the unit-lattice gradient of the block average of $\lambda$:
\[
\bQ(A+d\lambda)=\bQ A+\dl Q_0\lambda .
\]
In particular, the gauge transformations that preserve all averages are those with $Q_0\lambda$ constant. Modulo constants, these are those with $Q_0\lambda=0$. This is the residual gauge freedom of the variational problem, cf.\ \cite[(1.13), (1.20)]{BalPropI}.

\subsection{Restriction of product convolutions}

The averaging operators restrict a smoothed field to the unit lattice. The following estimate is used whenever we need to bound such restrictions in $\ell^2$.

\begin{lemma}[Restriction lemma]\label{lem:restriction}
Let $\cK=K_1\cdots K_d$ be a product convolution, where $K_\mu$ has a kernel bounded by $M_\mu$ and vanishing outside an interval of length $\ell_\mu\le6$. Then for every scalar field $f$,
\[
\|\iota\cK f\|_{\ell^2(\T)}\le\prod_{\mu=1}^dM_\mu(\ell_\mu+1)\ \|f\| .
\]
\end{lemma}

\begin{proof}
Write $(\cK f)(y)=\sum_t\eta^d\kappa(t)f(y+t)$ with $\kappa(t)=\prod_\mu\kappa_\mu(t_\mu)$. By Cauchy--Schwarz,
\[
|(\cK f)(y)|^2\le\Big(\sum_t\eta^d|\kappa(t)|\Big)\Big(\sum_t\eta^d|\kappa(t)|\,|f(y+t)|^2\Big).
\]
The first factor is at most $\prod_\mu M_\mu(\ell_\mu+\eta)$. Summing the second factor over $y\in\T$ gives $\sum_x\eta^d|f(x)|^2\sum_y|\kappa(x-y)|$. For fixed $x$ and $\mu$, at most $\ell_\mu+1$ integers $y_\mu$ satisfy $\kappa_\mu(x_\mu-y_\mu)\neq0$. Since $\ell_\mu+1\le7<n$, no point is counted twice on the torus. Hence $\sum_y|\kappa(x-y)|\le\prod_\mu M_\mu(\ell_\mu+1)$.
\end{proof}

When several convolutions act in the same direction, they are first merged into one directional convolution, using the composition rule above. For example, $Q_\sigma$ itself, $\iota\,a_\mu a_\nu b$, and $\iota\,b\,g_\mu$ satisfy the hypotheses with constants depending only on $d$. A shift $\tau_\mu$ composed with a convolution only moves the support of the kernel, so it does not change the bound.
